\scanpage{22}
A - 18 - \hfill n° 413

7. \uline{Noyaux formels}.

DÉFINITION 3.- \uline{Soit} $A$ \uline{une K-algèbre. On appelle carte formelle de} $A$ \uline{un couple} $\varphi=(M,\varphi^*)$ \uline{formé d'un module libre} $M$ \uline{sur} $K$ \uline{et d'un isomorphisme} $\varphi^*$ \uline{de} $\widehat P(M)$ \uline{sur l'algèbre} $A$. \uline{On dit encore que} $\varphi$ \uline{est à valeurs dans} $M$. \uline{On dit que deux cartes formelles} $\varphi_i$ $(i=1,2)$ \uline{de} $A$ \uline{à valeurs dans des modules libres} $M_i$ \uline{sont compatibles s'il existe une application formelle} $\psi_1$ \uline{de} $M_1$ \uline{dans} $M_2$ \uline{et une application formelle} $\psi_2$ \uline{de} $M_2$ \uline{dans} $M_1$ \uline{telles que} $\varphi_2^*=\psi_1^*\circ\varphi_1^*$ \uline{et} $\varphi_1^*=\psi_2^*\circ\varphi_2^*$.

Il est immédiat que la relation de compatibilité entre cartes formelles d'une même algèbre $A$ est une relation d'équivalence. Cela nous permet de poser la :

DÉFINITION 4.- \uline{On appelle noyau formel} $\underline V$ (\uline{sur l'anneau} $K$) \uline{la donnée d'une K-algèbre notée} $\widehat P(\underline V)$ \uline{et d'une classe d'équivalence de cartes formelles de} $\widehat P(\underline V)$.

Naturellement la donnée d'\uline{une} carte formelle d'une algèbre $A$ définit un noyau formel.

DÉFINITION 5.- \uline{Soient} $\underline V$ \uline{et} $\underline V'$ \uline{deux noyaux formels. Un morphisme} $f$ \uline{de} $\underline V$ \uline{dans} $\underline V'$ \uline{est la donnée d'un homomorphisme} $f^*$ \uline{de l'algèbre} $\widehat P(\underline V')$ \uline{dans l'algèbre} $\widehat P(\underline V)$ \uline{tel que pour toute carte formelle} $\varphi=(M,\varphi^*)$ \uline{de} $\underline V$ \uline{et toute carte formelle} $\psi=(N,\psi^*)$ \uline{de} $\underline V'$, \uline{l'homomorphisme} $\varphi^{*-1}\circ f^*\circ\psi^*$ \uline{de} $\widehat P(N)$ \uline{dans} $\widehat P(M)$ \uline{soit égal à l'homomorphisme} $\underline g^*$ \uline{défini par un élément} $\underline g\in\widehat P_+(M;N)$.

Il est immédiat qu'il suffit que le critère de la définition 5 soit vérifié pour un couple $(\varphi,\psi)$ de cartes formelles. Par ailleurs on vérifie aisément que les morphismes ainsi définis ne manquent d'aucune des mirifiques propriétés fonctorielles nécessaires pour avoir droit au soleil. On peut donc définir l'\uline{isomorphisme} de deux noyaux formels, ce que nous laisserons faire au prochain rédacteur, au lecteur, au Pape ou au Général.

\margin{\mspage{21}}
